%%%%%%%%%%%%%%%%%
% Andrea Adelfio - CV in AltaCV Format
%%%%%%%%%%%%%%%%%

\documentclass[9pt,a4paper,withhyper,normalphoto]{altacv}

\geometry{left=1.25cm,right=1.25cm,top=1.5cm,bottom=1.5cm,columnsep=0.6cm}
\usepackage{paracol}
\usepackage{ragged2e}

% Font
\renewcommand{\familydefault}{\sfdefault}
\usepackage{helvet}

% Europass-like colors
\definecolor{europassblue}{HTML}{003399}
\definecolor{SlateGrey}{HTML}{2E2E2E}
\definecolor{LightGrey}{HTML}{666666}

\colorlet{name}{black}
\colorlet{tagline}{europassblue}
\colorlet{heading}{europassblue}
\colorlet{headingrule}{europassblue}
\colorlet{subheading}{europassblue}
\colorlet{accent}{europassblue}
\colorlet{emphasis}{SlateGrey}
\colorlet{body}{LightGrey}

\renewcommand{\namefont}{\Huge\bfseries}
\renewcommand{\personalinfofont}{\footnotesize}
\renewcommand{\cvsectionfont}{\LARGE\bfseries}
\renewcommand{\cvsubsectionfont}{\large\bfseries}

\begin{document}

\name{Andrea Adelfio}
\tagline{Data Scientist | ML Engineer | Software Engineer}
\photoL{3.2cm}{picture}

\personalinfo{
  \email{and.adelfio@gmail.com}
  \phone{+39 389 162 5750}
  \location{Munich, Germany | \textbf{Open to EU Relocation}}
  \linkedin{andrea-adelfio-8544211b4}
  \github{andreaadelfio}
}

\makecvheader

\columnratio{0.28}
\begin{paracol}{2}


\cvsection{Tech Skills}

\cvsubsection{Programming}
\cvtag{Python (Expert)}\cvtag{SQL (Expert)}
\cvtag{Java (Proficient)}\cvtag{Bash}
\cvtag{JavaScript}\cvtag{C/C++}\cvtag{\LaTeX}

% \divider\smallskip

\cvsubsection{Machine Learning \& AI}
\cvtag{Deep Learning}\cvtag{Bayesian Neural Networks} \\
\cvtag{Anomaly Detection}\cvtag{Time Series Analysis} \\
\cvtag{TensorFlow, PyTorch, Keras}

% \divider\smallskip

\cvsubsection{Data Engineering \& DevOps}
\cvtag{Docker}\cvtag{Kubernetes}\cvtag{Git}\cvtag{CI/CD}\cvtag{HPC}\cvtag{Data Lakes}
\cvtag{Atlassian (Jira, Confluence...)}

\cvsection{Education}

\cvevent{M.Sc. in Physics (EQF 7)}{University of Trieste}{2018 -- 2020}{}
Thesis: \emph{Gravitational Wave Counterparts Detection with CTA}.  
Final grade: 101/110.

\divider

\cvevent{B.Sc. in Physics (EQF 6)}{University of Cagliari}{2015 -- 2018}{}
Final grade: 97/110.

\cvsection{Languages}
-\textbf{Italian} (Mother tongue) 

-\textbf{English} (C2 – Cambridge C2 Proficiency, score 202)  

-\textbf{German} (B1)

\cvsection{Certifications \& Courses}
-\textbf{School on Open Science Cloud} (kw: MLOps, Docker, Kubernetes)

-\textbf{Container Fundamentals for Scientific Research} (kw: Docker, Kubernetes)

-\textbf{High Performance Computing} (kw: HPC, c++, openmp)

-\textbf{Uncertainty Quantification for Machine Learning}

-\textbf{FPGA Programming}
\switchcolumn


\cvsection{Professional Profile}
\begin{quote}
\begin{justify}

Profile with five years of dual experience as Data Scientist, with a background in Artificial Intelligence and Machine Learning applied to data and time series analysis, and as Software Developer with expertise in Agile methodologies (Scrum) and CI/CD practices.

Ability to work across research and industrial contexts, with direct involvement in machine learning model development, software engineering and data infrastructures, in domains from the banking sector to satellite monitoring.
\end{justify}

\end{quote}

\cvsection{Work Experience}

\cvevent{Data Scientist \& ML Engineer}
{Italian Research Center on High Performance Computing, Big Data and Quantum Computing (ICSC) \& National Institute of Nuclear Physics (INFN)}
{Sep 2023 -- Dec 2025}
{Perugia, Italy}

\begin{itemize}
\item \textbf{ML Software Development:} Design, development and maintenance of \href{https://github.com/andreaadelfio/TSLies}{TSLies (link)}, an open-source Python framework for neural-network-based anomaly detection (FFNN, RNN, Bayesian Neural Networks) on multivariate time-series data (TensorFlow and TensorFlow-Probability -based pipeline). Applied to Fermi (NASA) satellite data.
\item \textbf{Mentoring \& Technical Leadership:} Supervision and mentoring of junior researchers and Master’s students, with guidance on software engineering best practices and code optimization. Contribution to two Master Classes for programming laboratories.
\item \textbf{Industrial Impact:} Implementation of the ATS framework for financial anomaly monitoring for Intesa Sanpaolo, bridging research and industry requirements.
\item \textbf{Data Infrastructure:} Design of scalable Data Lakes and storage architectures for the Italian Computing Center infrastructure, ensuring reliability and reproducibility for large-scale datasets.
\end{itemize}

\divider

\cvevent{Software Engineer}
{Accenture Technology Solutions}
{May 2021 -- Aug 2023}
{Milan, Italy}

\begin{itemize}
\item \textbf{Full-Stack Engineering:} Development and maintenance of mission-critical enterprise software using Java, SQL and JavaScript.
\item \textbf{Production Deployment:} Management of end-to-end software lifecycles, including version control, unit testing and deployment to production environments.
\item \textbf{Agile \& DevOps:} Work within Scrum teams, contributing to CI/CD pipelines and workflow automation.
\item \textbf{Consultancy:} Delivery of scalable technical solutions to international clients.
\end{itemize}

\cvsection{Key Projects}

\cvachievement{\faCode}{TSLies – Time Series anomaLies}
{Author and maintainer of an open-source Python package for anomaly detection in time series data based on neural networks and statistical learning methods.}

\divider

\cvachievement{\faChartLine}{Satellite Anomaly Detection}
{Design and implementation of modular ML pipelines for transient and anomaly detection on large-scale satellite time series data.}

\divider

\cvachievement{\faDatabase}{Data Infrastructure}
{Development of data-access, preprocessing pipelines and storage solutions for large datasets with reproducible workflows.}

\end{paracol}




\cvsection{Publications}

\begin{itemize}
\item “\textbf{Advances on TSLies: Time Series AnomaLIES, an Efficient Ma-
chine Learning Based Anomaly Detection Framework}”. Andrea Adelfio et al. in: Astronomy and Computing
(2025, under review).
\item “\textbf{Scalable and Interoperable Data Management in the Spoke 3
Big Data Infrastructure}”. Giacomo Coran et al. in: Astronomy and Computing (2025, under review)
\item \href{https://doi.org/10.1109/PDP66500.2025.00062}{“\textbf{Anomaly Detection with Machine Learning on Time Series: Unveiling Lost Transients Data}”. Andrea Adelfio et al., in: Proceedings of the Euromicro International Conference on Parallel, Distributed, and Network-Based Processing (PDP). 2025.}
\item \href{https://doi.org/10.1109/PDP66500.2025.00072}{“\textbf{Spoke 3 Big Data Management, Storage, and Archive Infrastructure}”. Giacomo Coran et al. in: Proceedings of the Euromicro International Conference on Parallel, Distributed, and Network-Based Processing (PDP). 2025.}
\item \href{https://doi.org/10.22323/1.395.0998} {“\textbf{Searching for Very-High-Energy Electromagnetic Counterparts to Gravitational-Wave Events with the Cherenkov Telescope Array}”. Barbara Patricelli et al. in: Proceedings of the 37th International Cosmic Ray Conference. 2021.}
\end{itemize}
\end{document}